\section*{Chapitre \ref{ch:dv:the-binomial-model} --- Le modèle binomial}

\begin{solution}{exo:dv:the-binomial-model:1}
$\Delta=(10-0)/(60-40)=0.5$ action~; le dépôt $b$ vérifie $0.5\times40+b=0$, donc
$b=-20$ (un emprunt). Prix $0.5\times50-20=5$~; $p=(1-0.8)/(1.2-0.8)=0.5$.
\end{solution}

\begin{solution}{exo:dv:the-binomial-model:2}
$u=e^{0.2/\sqrt{12}}=1.0594$, $d=1/u=0.9439$,
$p=(e^{0.03/12}-d)/(u-d)=0.5072$.
\end{solution}

\begin{solution}{exo:dv:the-binomial-model:3}
$C-P=10.3603-4.5925=5.7678$ et $S-KR^{-3}=100-100/1.02^3=5.7678$.
\end{solution}

\begin{solution}{exo:dv:the-binomial-model:4}
Le call est construit avec des actions et des liquidités, nœud par nœud, pour 10.3603, et verse le
gain du call sur tous les chemins, quelle que soit la probabilité. Si le call se négociait au-dessus de ce niveau,
la vendre et acheter le portefeuille constituerait un arbitrage~; en dessous, ce serait l'inverse.
Le 0.8 du collègue modifie le gain espéré du call, et de même le
rendement espéré de l'action~: il modifie les primes de risque, et non le prix relatif à l'action.
\end{solution}

\begin{solution}{exo:dv:the-binomial-model:5}
Il paie aux nœuds 133.1 et 108.9~: $(0.6^3+3\times0.6^2\times0.4)/1.02^3=0.648/1.0612=0.6106$.
\end{solution}

\begin{solution}{exo:dv:the-binomial-model:6}
Cox--Ross--Rubinstein~: $u=1.0653$, $d=0.9387$, $p=0.5238$. Jarrow--Rudd~:
$u=1.0685$, $d=0.9415$, $p=0.5000$ à quatre décimales près. La première place la dérive
dans la probabilité, la seconde dans les facteurs.
\end{solution}

\begin{solution}{exo:dv:the-binomial-model:7}
Cox--Ross--Rubinstein 5.5909 (0.0174 de trop)~; Leisen--Reimer 5.5735 (erreur
$-0.00003$).
\end{solution}

\begin{solution}{exo:dv:the-binomial-model:8}
Une probabilité de 1.125 signifie que la liquidité croît plus vite que l'action ne peut le faire dans
l'un ou l'autre état ($R=1.06>u=1.05$)~: vendre l'action à découvert et déposer le produit de la vente rapporte un
profit sans risque. La «~calibration~» décrit un arbitrage, et les prix calculés
avec elle n'ont aucun sens. Le modèle exige $d<R<u$.
\end{solution}

\begin{solution}{pb:dv:the-binomial-model:1}
\textbf{1.} $p=(1.02-0.9)/0.2=0.6$.
\textbf{2.} 133.1, 108.9, 89.1, 72.9.
\textbf{3.} 10.3603.
\textbf{4.} 0.6240 action et un dépôt de $-52.04$.
\textbf{5.} 4.5925~; $C-P=5.7678=100-100/1.02^3$.
\textbf{6.} À 90 après un pas (10 contre 9.196 si l'on conserve l'option) et à 81 après deux (19
contre 17.039).
\textbf{7.} 4.9076.
\textbf{8.} 0.3151.
\textbf{9.} Non~: 10.3603 dans les deux cas, car exercer un call par anticipation renonce à l'intérêt
sur le prix d'exercice et à la valeur résiduelle de l'option.
\textbf{10.} L'exercice rapporte 10 immédiatement~; détenir l'option vaut 9.196, car le gain
peut au mieux valoir 19 deux pas plus tard et est souvent moindre~: l'intérêt sur le
prix d'exercice l'emporte sur la valeur optionnelle restante.
\textbf{11.} 5.5735.
\textbf{12.} $-0.0200$ à 100 pas, $+0.0174$ à 101.
\textbf{13.} $-0.00003$.
\textbf{14.} 6.0902~; prime d'exercice anticipé de 0.5167 par rapport à la valeur européenne donnée par la formule.
\textbf{15.} Avec le prix d'exercice égal au spot, un $n$ pair place un nœud sur le coude et un $n$ impair
l'encadre~: le gain discrétisé est alternativement trop faible et trop élevé près
du coude.
\textbf{16.} Cox--Ross--Rubinstein 199 pas~; Leisen--Reimer 7 (tout impair $n\ge7$
est à moins d'un centime près).
\textbf{17.} En $n^2$~: 2\,001 pas coûtent environ 400 fois le travail de 100. Pour
10\,000 options, le desk utilise l'arbre précis pour les petits $n$, ou une grille
où le même travail est partagé entre les prix d'exercice.
\textbf{18.} Deux mouvements possibles par pas, couverture continue à la limite, sans
coûts, une volatilité connue~: il valorise par réplication dans un marché qui n'est pas
complet.
\textbf{19.} 0.3151 au tableau~; 199 pas.
\textbf{20.} C'est le poids qui fait que les actions et les liquidités reproduisent mutuellement
leurs prix, fixé par $u$, $d$ et $R$, quelles que soient les anticipations sur l'avenir.
\end{solution}

\begin{solution}{iq:dv:the-binomial-model:1}
$\Delta=20/40=0.5$ action, emprunter 40 (de sorte que $0.5\times80-40=0$)~: le call coûte
$50-40=10$~; $p=0.5$.
\iqlookfor{réplication, et non une valeur espérée sous une probabilité supposée.}
\end{solution}

\begin{solution}{iq:dv:the-binomial-model:2}
Parce que l'option est répliquée par des actions et des liquidités sur chaque chemin~; son prix est
le coût du \omterm{def:dv:no-arbitrage-and-the-fundamental-theorems:claim}{portefeuille de réplication}. Toute modification de la probabilité change le
rendement espéré de l'action aussi, et le prix relatif reste inchangé.
\iqlookfor{réplication~; la probabilité comme poids dérivé.}
\end{solution}

\begin{solution}{iq:dv:the-binomial-model:3}
Le coude du gain se situe à un endroit différent par rapport aux nœuds terminaux
selon que $n$ est pair ou impair, si bien que le gain échantillonné alterne entre trop élevé et trop
faible. Remède~: Leisen--Reimer (centrer l'arbre sur le prix d'exercice), moyenner $n$ et $n+1$,
ou lisser le gain sur le dernier pas (chapitre 22).
\iqlookfor{le coude comme cause.}
\end{solution}

\begin{solution}{iq:dv:the-binomial-model:4}
Un put profond dans la monnaie vaut presque $K-S$ à l'échéance~; l'exercice immédiat rapporte
l'intérêt sur $K$ et fait perdre peu de valeur optionnelle. Pour un call, exercer verse $K$
par anticipation et renonce à l'assurance de la valeur optionnelle, alors que $C\ge S-Ke^{-r(T-t)}>
S-K$~: il vaut toujours mieux le vendre que de l'exercer.
\iqlookfor{l'intérêt sur le prix d'exercice face à la valeur optionnelle.}
\end{solution}

\begin{solution}{iq:dv:the-binomial-model:5}
$(n+1)(n+2)/2$ nœuds~; la récurrence rétrograde n'exige qu'une seule tranche temporelle à la fois,
$n+1$ valeurs, d'où une mémoire en $O(n)$ et un temps en $O(n^2)$.
\iqlookfor{recombinaison et tableau à une seule tranche.}
\end{solution}

\begin{solution}{iq:dv:the-binomial-model:6}
On les lit sur les nœuds après un et deux pas~: $\Delta\approx(V_1^1-V_1^0)/
(Su-Sd)$ et le gamma à partir des trois nœuds du pas 2, centrés sur $S$ dans un
\omterm{def:dv:the-binomial-model:crr}{arbre de Cox--Ross--Rubinstein} ($Su^2$, $S$, $Sd^2$). Mieux vaut commencer l'arbre deux pas
avant aujourd'hui, de sorte que les trois nœuds du jour soient $Su^2$, $S$, $Sd^2$ et que les
différences soient centrées.
\iqlookfor{grecques calculées à partir des nœuds de l'arbre lui-même, centrés.}
\end{solution}
